\section{需求分析}
平台围绕医生、患者与管理员三类角色展开。医生侧关注病历整合效率、风险识别准确性与告警可解释性；患者侧关注提问便捷性、建议可执行性与随访连续性；管理员侧关注模板配置、模型版本管理与审计追踪能力。结合开题阶段调研，系统需满足以下核心需求：
\begin{itemize}
  \item \textbf{数据治理需求}：支持文档导入、结构化抽取、字段标准化与质量校验；
  \item \textbf{智能交互需求}：提供基于知识库的问答服务，支持多轮上下文与溯源展示；
  \item \textbf{风险预警需求}：基于患者特征生成风险评分并支持阈值化告警；
  \item \textbf{运维治理需求}：支持Prompt版本管理、日志留痕与模型效果回溯。
\end{itemize}

\section{总体架构}
系统采用前后端分离与微服务架构，将业务能力按“接入层—服务层—AI与数据层”进行拆分。接入层提供医生端Web与患者端交互界面；服务层承担认证授权、病历管理、问答编排、风险预测与模板管理；AI与数据层负责向量检索、结构化存储及模型推理。

\begin{figure}[htbp]
  \centering
  \includegraphics[width=0.92\textwidth]{system-architecture.png}
  \caption{平台总体架构图}
\end{figure}

\section{关键业务流程}
在典型业务流程中，患者数据经OCR或接口接入后进入结构化处理阶段，随后由风险服务完成特征提取与评分计算。对于高风险样本，系统触发告警并向医生提供可解释的关键因子提示；对于中低风险样本，系统生成个性化健康建议与复诊提醒。问答服务在全流程中提供即时解释能力，支持患者在院外阶段的连续咨询。该流程实现了“数据入库—风险识别—干预建议—结果回写”的闭环管理。